\documentclass[11pt]{article}
\usepackage[margin=1in]{geometry}
\usepackage{amsmath, amssymb}
\usepackage{xcolor}

\newcommand{\E}{\mathbb{E}}

\title{Why the Analytical Poisson Entropy Formula Avoids Truncation Errors}
\author{Explanation for Monte Carlo vs Analytical Comparison}
\date{}

\begin{document}
\maketitle

\section{The Standard Entropy Formula}

For a Poisson distribution with rate $f$, the probability mass function is:
\begin{equation}
p(r \mid f) = \frac{f^r e^{-f}}{r!}, \quad r = 0, 1, 2, \ldots
\end{equation}

The entropy is defined as:
\begin{equation}
H(f) = -\sum_{r=0}^{\infty} p(r \mid f) \log p(r \mid f)
\label{eq:entropy_def}
\end{equation}

\section{Expanding the Log-Probability}

First, we expand the log of the Poisson probability:
\begin{align}
\log p(r \mid f) &= \log\left( \frac{f^r e^{-f}}{r!} \right) \\
&= r \log(f) - f - \log(r!)
\end{align}

\section{Substituting into the Entropy Formula}

Substituting this into \eqref{eq:entropy_def}:
\begin{align}
H(f) &= -\sum_{r=0}^{\infty} p(r) \left[ r \log(f) - f - \log(r!) \right] \\
&= -\log(f) \sum_{r=0}^{\infty} r \cdot p(r) + f \sum_{r=0}^{\infty} p(r) + \sum_{r=0}^{\infty} p(r) \log(r!)
\label{eq:entropy_expanded}
\end{align}

\section{Using Known Properties of the Poisson Distribution}

We know two fundamental properties of any probability distribution:

\begin{enumerate}
\item \textbf{Normalization:} $\displaystyle\sum_{r=0}^{\infty} p(r) = 1$

\item \textbf{Poisson mean:} $\displaystyle\sum_{r=0}^{\infty} r \cdot p(r) = f$
\end{enumerate}

These are \textbf{exact identities} that hold for the infinite sum. We don't need to compute them numerically.

\section{The Decomposed Entropy Formula}

Substituting these known values into \eqref{eq:entropy_expanded}:
\begin{align}
H(f) &= -\log(f) \cdot f + f \cdot 1 + \sum_{r=0}^{\infty} p(r) \log(r!) \\
&= f - f\log(f) + \E_r[\log(r!)] \\
&= \boxed{f(1 - \log f) + \E_r[\log(r!)]}
\label{eq:decomposed}
\end{align}

\textbf{Key insight:} The first term $f(1 - \log f)$ is \textbf{exact} --- it uses the analytical values 1 and $f$, not truncated sums.

\section{What Happens with Truncation}

In practice, we truncate the sum at some $R_{\max}$ (e.g., 100). Consider what the \textbf{naive approach} computes:
\begin{equation}
H_{\text{trunc}}(f) = -\sum_{r=0}^{R_{\max}-1} p(r) \log p(r)
\end{equation}

Expanding as before:
\begin{align}
H_{\text{trunc}}(f) &= -\log(f) \underbrace{\sum_{r < R_{\max}} r \cdot p(r)}_{M_{\text{trunc}}} + f \underbrace{\sum_{r < R_{\max}} p(r)}_{S_{\text{trunc}}} + \sum_{r < R_{\max}} p(r) \log(r!)
\end{align}

\subsection{The Problem: Truncated Sums Are Wrong}

For large $f$ (e.g., $f = 150$), significant probability mass lies beyond $r = 100$.

\begin{center}
\begin{tabular}{|c|c|c|}
\hline
\textbf{Quantity} & \textbf{True Value} & \textbf{Truncated Value ($R_{\max}=100$)} \\
\hline
$\sum p(r)$ & 1 & $S_{\text{trunc}} \approx 0.0001$ \\
$\sum r \cdot p(r)$ & 150 & $M_{\text{trunc}} \approx 0.01$ \\
\hline
\end{tabular}
\end{center}

The truncated formula gives:
\begin{equation}
H_{\text{trunc}} = -\log(f) \cdot M_{\text{trunc}} + f \cdot S_{\text{trunc}} + \ldots \neq f(1 - \log f) + \ldots
\end{equation}

\textcolor{red}{\textbf{This is completely wrong!}} The truncation corrupts the $f(1-\log f)$ term.

\section{The Correct Approach}

Use the decomposed formula \eqref{eq:decomposed}:
\begin{equation}
H(f) = \underbrace{f(1 - \log f)}_{\text{compute exactly}} + \underbrace{\sum_{r=0}^{R_{\max}-1} p(r) \log(r!)}_{\text{truncated sum (unavoidable)}}
\end{equation}

\begin{itemize}
\item The first term uses $f$ directly --- \textbf{no summation needed, no truncation error}.
\item The second term still has truncation, but this is the \textbf{same truncation} used in the analytical formula.
\end{itemize}

\section{Application to the Monte Carlo Test}

\subsection{What \texttt{nd\_utility} Does (Analytical)}

The expected noise entropy is:
\begin{equation}
\E_{\lambda}[H(\exp(\lambda))] = \E[f(1 - \log f)] + \E[\E_r[\log(r!)]]
\end{equation}

where $\lambda \sim \mathcal{N}(\mu, \sigma^2)$ and $f = e^\lambda$.

Using the moment generating function of the normal distribution:
\begin{align}
\E[e^\lambda] &= e^{\mu + \sigma^2/2} \\
\E[\lambda \cdot e^\lambda] &= (\mu + \sigma^2) e^{\mu + \sigma^2/2}
\end{align}

Therefore:
\begin{align}
\E[f(1 - \log f)] &= \E[e^\lambda(1 - \lambda)] \\
&= \E[e^\lambda] - \E[\lambda \cdot e^\lambda] \\
&= e^{\mu + \sigma^2/2} - (\mu + \sigma^2) e^{\mu + \sigma^2/2} \\
&= \boxed{e^{\mu + \sigma^2/2}(1 - \mu - \sigma^2)}
\end{align}

This is \textbf{exact} --- computed in closed form, no infinite sums.

\subsection{What the Original Monte Carlo Did (Wrong)}

For each sample $\lambda_j$:
\begin{equation}
H_j = -\sum_{r=0}^{99} p(r \mid e^{\lambda_j}) \log p(r \mid e^{\lambda_j})
\end{equation}

When $e^{\lambda_j}$ is large, this gives the \textbf{wrong} entropy due to truncation.

\subsection{What the Fixed Monte Carlo Should Do}

For each sample $\lambda_j$, compute:
\begin{equation}
H_j = \underbrace{e^{\lambda_j}(1 - \lambda_j)}_{\text{exact}} + \underbrace{\sum_{r=0}^{99} p(r \mid e^{\lambda_j}) \log(r!)}_{\text{truncated, same as analytical}}
\end{equation}

Then average: $H_{\text{cond}} = \frac{1}{N}\sum_j H_j$.

This should match the analytical formula.

\section{Summary}

\begin{center}
\begin{tabular}{|l|c|c|}
\hline
\textbf{Method} & \textbf{$f(1-\log f)$ term} & \textbf{$\E[\log(r!)]$ term} \\
\hline
Analytical (\texttt{nd\_utility}) & Exact (closed form) & Truncated sum \\
Original Monte Carlo & \textcolor{red}{Wrong (truncated)} & Truncated sum \\
Fixed Monte Carlo & Exact ($e^\lambda(1-\lambda)$) & Truncated sum \\
\hline
\end{tabular}
\end{center}

The original Monte Carlo had 8\% error because the truncation corrupted the $f(1-\log f)$ term. The fix computes this term exactly.

\end{document}
